\chapter{Background}
\label{chap:background}


\section{Sviluppo di applicazioni full stack monolinguaggio}
\label{sec:background-fullstack}
Prima di addentrarci nella descrizione del background del progetto è importante definire in maniera corretta cosa sono le applicazioni \textbf{full stack} e le differenze che emergono tra lo sviluppo di esse e quello di applicazioni più tradizionali.
Le componenti delle applicazioni vengono divise in due macrocategorie
\begin{itemize}
    \item \textbf{Frontend}: in questa sezione dell'applicativo sono gestiti tutti gli aspetti
    inerenti all'interfaccia utente: sia quelli legati all'aspetto grafico che quelli legati alla 
    logica di interazione con essa.
    \item \textbf{Backend}: in questa sezione invece sono gestiti tutti gli aspetti legati\\
    all'esecuzione delle funzionalità dell'applicazione, inclusa l'elaborazione dei dati, la loro
    memorizzazione e la comunicazione con altre applicazioni se necessario. 
\end{itemize}
Le applicazioni \textit{full stack} dunque, comprendono entrambe le componenti.
Lo sviluppo di questo tipo di applicazioni può includere l'utilizzo di molteplici linguaggi in base alle necessità, impiegando linguaggi progettati per lo sviluppo del backend e altri per il frontend.
Tuttavia questa distinzione al giorno d'oggi viene gradualmente smussata, infatti diversi linguaggi inizialmente progettati per l'utilizzo in una sola area, vengono poi estesi anche all'altra.
L'esempio più rilevante e dominante al giorno d'oggi è \textbf{JavaScript}, il quale inizialmente nasce come linguaggio per i browser, quindi dedicato al \textit{frontend}, per poi evolversi oltre i confini del browser permettendone l'esecuzione anche sui server. 
Questo permette dunque la creazione di applicativi \textit{full stack} monolinguaggio, una tendenza sempre più diffusa.
Tra i vantaggi che essa comporta vi sono l'utilizzo di un solo linguaggio per tutto lo stack, il quale porta a una minore curva di apprendimento e una maggiore efficienza per quanto riguarda l'impiego di programmatori.
Uno dei più grandi limiti dello sviluppo \textit{full stack} è l'interazione con librerie e API
native, infatti linguaggi come JavaScript sono noti per avere meccanismi limitati o comunque ostici per interagire con il codice nativo.
In seguito verranno discusse in dettaglio le soluzioni esistenti per la programmazione di
tale tipo di applicativi e i loro limiti, per poi analizzare il problema della portabilità e
dell'integrazione con le librerie native.

\subsection{Panoramica delle soluzioni esistenti e loro limiti}
\label{subsec:background-fullstack-soluzioni}
Come anticipato in precedenza, esistono diversi linguaggi impiegati per la programmazione di app \textit{full stack} monolinguaggio.
Il più popolare attualmente risulta essere JavaScript \cite{stackoverflow2025}, questo per via dello sviluppo della tecnologia \textbf{Node.js}; infatti questo permette di portare l'ambiente di esecuzione di JavaScript fuori dai browser web.
Inoltre anche la creazione di TypeScript, una diretta evoluzione di JavaScript, contribuì alla crescita di questo tipo di linguaggi.
Tutti questi fattori rendono JavaScript una delle scelte più gettonate per la creazione di software \textit{full stack}, impiegando diversi stack come Next.js, framework sviluppato da Vercel, e lo stack MEAN, il quale consiste in Mongo DB, Express.js e Angular come framework per il backend e frontend rispettivamente, il tutto legato da Node.js.
Nonostante il grande supporto in termini di librerie e framework per questo tipo di tecnologia, l'integrazione con le librerie native risulta comunque ardua per via delle peculiarità del linguaggio.
Esistono anche soluzioni basate su linguaggi eseguiti su \textit{Java Virtual Machine}, come Java e Kotlin. Per questi esistono framework maturi sia per il backend, come Ktor, che per il frontend, come Vaadin. Tuttavia queste soluzioni rimangono legate all'ecosistema JVM, il che pone vincoli sulla portabilità verso alcune piattaforme target.
Nonostante le differenze tra queste soluzioni, emerge un limite comune: l'integrazione con librerie native risulta in tutti i casi problematica, come verrà analizzato nella sezione seguente.


%Ad esempio per poter interagire con una libreria C++, Node.js mette a disposizione \textbf{N-API}, una libreria di C++ creata appositamente per fare da ponte tra la libreria C++ che si vuole utilizzare e l'ambiente di esecuzione Node.js \cite{nodejs_napi}.
%Questo però contraddice direttamente la filosofia dello sviluppo \textit{full stack} monolinguaggio, rendendo necessario per lo sviluppatore conoscere linguaggi come C++ per poter scrivere il collante tra la libreria nativa e JavaScript.
%Altre opzioni per lo sviluppo \textit{full stack} le troviamo ad esempio anche con Python, che al contrario di Node.js presenta varie soluzioni per interagire con il nativo, approccio però che aumenta la frammentazione e spesso coinvolge la scrittura di codice in altri linguaggi per la compilazione nativa, oppure scrittura di codice per tradurre le API a Python.


\subsection{La portabilità e integrazione con librerie native}
\label{subsec:background-fullstack-gap}

\section{Kotlin e framework multitarget come risposta al gap}
\label{sec:background-multitarget}

\subsection{Kotlin Multiplatform}
\label{subsec:background-kmp}

\subsection{Interoperabilità con C tramite Cinterop}
\label{subsec:background-kmp-cinterop}

\section{Lo standard FMI e il formato FMU}
\label{sec:background-fmi}

\subsection{Functional Mock-up Interface (FMI 2.0)}
\label{subsec:background-fmi-standard}

\subsection{Struttura di un file FMU e modalità operative}
\label{subsec:background-fmu-struttura}